% Part: conditional-logics
% Chapter: minimal-change-semantics
% Section: true-false

\documentclass[../../../include/open-logic-section]{subfiles}

\begin{document}

\olfileid{con}{min}{tf}

\olsection{صدق الشرطيات المخالفة للواقع وكذبها}

يكون الشرط المخالف للواقع~$!A \cif !B$ صادقًا (على نحو غير فارغ) إذا
كانت جميع أقرب عوالم~$!A$ عوالم~$!B$، كما يبيّن
\olref{fig:true}.
\begin{figure}
\begin{center}
\begin{tikzpicture}[scale=.7]\small
  \spheresystem{5}
  \draw (0,0) node {$w$};
  \propositionintersect{0}{4}{40}{3.5}
  {\tikzset{proposition/.append={smooth,tension=1.4}}
  \proposition[shift={(1.6,-1)}]{90}{1}{30}{4}}
  \path (1.5,3) node[below] {$\formula{B}$};
  \path (3,0) node {$\formula{A}$};
\end{tikzpicture}
\caption{شرط مخالف للواقع صادق على نحو غير فارغ}
\ollabel{fig:true}
\end{center}
\end{figure}
ويكون الشرط المخالف للواقع صادقًا أيضًا عند~$w$ إذا لم يحتو نظام الكرات
المحيطة بـ$w$ أي كرة تقبل~$!A$ أصلًا. وفي هذه الحالة يكون صادقًا
\emph{على نحو فارغ} (انظر~\olref{fig:vacuous}).
\begin{figure}
\begin{center}
\begin{tikzpicture}[scale=.7]\small
  \spheresystem{5}
  \draw (0,0) node {$w$};
  \proposition{0}{6.5}{40}{3.5}
  {\tikzset{proposition/.append={smooth,tension=1.4}}
  \proposition[shift={(1.2,-1)}]{90}{1}{30}{4}}
  \path (1.5,3) node[below] {$\formula{B}$};
  \spherepos{0}{7}{node {$\formula{A}$}}
\end{tikzpicture}
\caption{شرط مخالف للواقع صادق على نحو فارغ}
\ollabel{fig:vacuous}
\end{center}
\end{figure}

وقد يكون كاذبًا بطريقتين. إحداهما أن تكون بعض أقرب عوالم~$!A$،
لا كلها، عوالم~$!B$. وفي هذه الحالة تكون~$!A \cif \lnot !B$ كاذبة
أيضًا (انظر~\olref{fig:false}).
\begin{figure}
\begin{center}
\begin{tikzpicture}[scale=.7]\small
  \spheresystem{5}
  \draw (0,0) node {$w$};
  \propositionintersect{0}{4}{40}{3.5}
  {\tikzset{proposition/.append={smooth,tension=1.4}}
  \proposition[shift={(1.6,0)}]{90}{1}{30}{3}}
  \path (1.5,3) node[below] {$\formula{B}$};
  \path (3,0) node {$\formula{A}$};
\end{tikzpicture}
\caption{شرط مخالف للواقع كاذب، ومقابله كاذب}
\ollabel{fig:false}
\end{center}
\end{figure}
أما إذا لم تتقاطع أقرب عوالم~$!A$ مع عوالم~$!B$ أصلًا، كانت
$!A \cif !B$ كاذبة. لكن جميع أقرب عوالم~$!A$ تكون في هذه الحالة
عوالم~$\lnot !B$، ولذلك تكون~$!A \cif \lnot !B$ صادقة
(انظر~\olref{fig:false-opposite}).
\begin{figure}
\begin{center}
\begin{tikzpicture}[scale=.7]\small
  \spheresystem{5}
  \draw (0,0) node {$w$};
  \propositionintersect{0}{4}{40}{3.5}
  {\tikzset{proposition/.append={smooth,tension=1.4}}
  \proposition[shift={(-1,0)}]{90}{1}{30}{3}}
  \path (-1,3) node[below] {$\formula{B}$};
  \path (3,0) node {$\formula{A}$};
\end{tikzpicture}
\caption{شرط مخالف للواقع كاذب، ومقابله صادق}
\ollabel{fig:false-opposite}
\end{center}
\end{figure}

بخلاف الشرط الصارم، يمكن أن تكون الشرطيات المخالفة للواقع عرضية.
انظر نموذج الكرات في~\olref{fig:contingent}. فجميع عوالم~$!A$ الأقرب
إلى~$u$ هي عوالم~$!B$، ولذلك~$\mSat{M}{!A \cif !B}[u]$.
لكن بعض عوالم~$!A$ الأقرب إلى~$v$ ليست عوالم~$!B$، ولذلك
$\mSat/{M}{!A \cif !B}[v]$.
\begin{figure}
\begin{center}
\begin{tikzpicture}[scale=.6]\small
  \spheresystem{7}
  \spheresystem[shift={(0:3.1)},dashed]{4}
  \propositionintersect{45}{4}{40}{4.5}
  \begin{scope}
    \clip \propositionplot{45}{4}{40}{4.5} ;
    \spherelayer[shift={(0:3.1)}]{3}
  \end{scope}
  \proposition[shift={(1.4,-.5)}]{90}{1}{35}{5}
  \path (0,0) node {$u$};
  \path (0:3.1) node {$v$};
  \spherepos{35}{9}{node {$\formula{A}$}}
  \spherepos{80}{9}{node {$\formula{B}$}}
\end{tikzpicture}
\end{center}
\caption{شرط مخالف للواقع عَرَضي}
\ollabel{fig:contingent}
\end{figure}
\end{document}
